%% P17 --- The Hessian, curvature and the Taylor expansion
%%
%% Every number in this program that the reader cannot do in their head comes
%% from code/p17_hessian_taylor.py and is pulled in with \val{}. The console
%% listing is a committed transcript written by that script.
%%
%% THE THESIS: the second-order model turns "the loss exploded" into
%% arithmetic. On the quadratic model each eigendirection is multiplied by
%% 1 - eta lambda every step, so bounded means |1 - eta lambda| < 1 in every
%% direction at once, which is eta < 2/lambda_max. A number, not a feeling.
%%
%% WHAT THIS PROGRAM IS OWED, read out of the files rather than remembered:
%%   P15  PUT THE READER ONE LINE FROM IT. It derives the factor per
%%        eigendirection, measures it as -0.80 and +0.91 on P10's bowl, and
%%        its rigour box says: "this program can say that the factor goes
%%        negative and that the walk crosses; it may not tell you where the
%%        boundary is." The boundary is here, and P15's own eta is GATED to
%%        sit under it.
%%   P10  built the bowl -- eigenvalues 20 and 1, the level ellipse, the words
%%        ravine and sharp minimum -- and said the inequality is P17's.
%%   F04  gives the sum and the sequence, which is what a Taylor expansion is.
%%   F11  gives the stationary point; F12 the derivatives it is built from.
%%   P16  priced a Jacobian. A Hessian is one dimension worse, and the count
%%        is gated against P16's own parameter figure.
%%
%% WHAT IT LEAVES ALONE, checked against tools/programs.json:
%%   matrix calculus, the softmax-cross-entropy gradient          -> P18
%%   convexity, one basin, what it promises                       -> P19
%%   the optimisers that live with the shape                      -> P20
%%   minibatch noise                                              -> P21
%%
%% SECTION 5 IS A DEMONSTRATION RATHER THAN A HEDGE. Sharpness is not a
%% property of the function a network computes: rescaling a parameter changes
%% the curvature while leaving every output identical, exactly, and the script
%% computes it rather than the prose arguing it.
%%
%% Twin: programs/pl/P17-hessian-taylor.tex. Frame for frame, macro for macro,
%% number for number; tools/parity.py compares them.

\program[The Hessian and curvature]{The Hessian, curvature and the Taylor
  expansion}
\label{prog:P17}
\index{Hessian}
\index{curvature}

\begin{outcomes}
\outcome{1--8}{say what the second-order Taylor model is and in what sense it
  is the best local one}
\outcome{9--15}{write the matrix of second derivatives, and say why
  Program~\ref{prog:P10} applies to it unchanged}
\outcome{16--23}{derive the largest step size a quadratic will tolerate, and
  say which direction sets it}
\outcome{24--30}{say what using the curvature directly would buy and what it
  would cost}
\outcome{31--38}{say precisely which claims about sharp and flat minima are
  about the function and which are about the coordinates}
\end{outcomes}

\begin{quiz}
\item You know a function's value, slope and curvature at a point. What is the
  best you can say about it nearby? \teachesat{3--4}
  \answerto{A parabola \dash{} the second-order Taylor model. Its error falls
  like the cube of the distance, where the tangent line's falls like the
  square, which is what \enquote{best} means here: an order, not a size.}
\item Program~\ref{prog:P15} found each direction multiplied by
  $1 - \eta\lambda$ every step. What does that make the largest usable $\eta$?
  \teachesat{16--17}
  \answerto{$2/\lambda$, because the walk stays bounded exactly when
  $\lvert 1 - \eta\lambda \rvert < 1$. With several directions the smallest of
  those bounds wins, so it is $2/\lambda_{\max}$.}
\item Which direction of a valley decides the step size, and which one has to
  live with it? \teachesat{18--19}
  \answerto{The steepest sets the cap, because it is the one that diverges
  first; the shallowest lives with it, and its own bound is larger by the
  condition number.}
\item Newton's method uses the curvature and lands on the minimum of a
  quadratic in one step. Why is it not used for training?
  \teachesat{26--27}
  \answerto{Because the Hessian has one entry per \emph{pair} of parameters
  and solving with it is worse still. At a real model size the matrix alone is
  larger than any storage that exists.}
\item Is \enquote{this minimum is sharper than that one} a fact about the
  trained model? \teachesat{31--33}
  \answerto{Not on its own. Rescaling a parameter changes the curvature while
  leaving every output identical, so a bare sharpness number is a fact about
  the coordinates. What survives reparameterisation has to be said carefully,
  and section 5 says which parts do.}
\end{quiz}

% ============================================================
\section{The best local model}
\label{sec:P17-taylor}
\index{Taylor expansion}

\begin{fr}
Program~\ref{prog:F11} approximated a curve near a point by its tangent line,
and Program~\ref{prog:P15} did the same in many dimensions with the gradient.
A tangent line is a \emph{model}: it agrees with the function at the point and
in slope, and disagrees everywhere else.

You have the value, the slope and now the rate at which the slope itself
changes. What would you build out of the three, and what shape would it be?
\dotline
\nextframe
\end{fr}

\begin{fr}
\begin{ansblock}
A parabola \dash{} value, plus slope times the distance, plus half the
curvature times the distance squared.
\end{ansblock}

\[ f(x_0 + h) \;\approx\; f(x_0) + f'(x_0)\,h + \tfrac{1}{2}f''(x_0)\,h^{2} \]

That is the \textbf{second-order Taylor model}, and the $\tfrac{1}{2}$ is not
a fudge: it is what makes the second derivative of the model equal the second
derivative of the function, which is the whole point of including the term.

\mermaidfig{p17-three-models}{Constant, linear, quadratic \dash{} one more
agreement each time.}{three local models}
\end{fr}

\begin{fr}
Now the claim that makes it worth having, and it is easy to state loosely.

It is tempting to say the parabola is \enquote{more accurate} than the tangent
line. That is true and it is not the useful statement. What is the useful one?
\nextframe
\end{fr}

\begin{fr}
\begin{ansblock}
Its error falls faster as you shrink the distance \dash{} like $h^{3}$ where
the line's falls like $h^{2}$. An \emph{order}, not a size.
\end{ansblock}

\textbf{That distinction is the difference between a claim about one step and a
claim about all of them.} \enquote{More accurate} is a fact about whatever $h$
you happened to try; \enquote{one order better} says what happens as you close
in, and it is checkable.

So the script checks it that way: it halves $h$ eight times and asserts the
linear error divides by about $2^{\val{p17.order.lin}}$ each time and the
quadratic error by about $2^{\val{p17.order.quad}}$. At
$h = \val{p17.taylor.h}$ the two errors are $\val{p17.err.lin}$ and
$\val{p17.err.quad}$, but those two figures are illustrations; the ratios are
the result.

\begin{note}
This is the same discipline Program~\ref{prog:P05} used for its spread and
Program~\ref{prog:P15} for its zig-zag angle: \textbf{assert the invariant, not
the observation.} Two error figures would have been a fact about one step size
and one function. The order survives both.
\end{note}
\end{fr}

\begin{fr}
One thing the model does not do, and saying so now saves trouble later.

The Taylor model is \emph{local}. Write down what that costs you, in terms of
the picture Program~\ref{prog:P15} drew.
\dotline
\nextframe
\end{fr}

\begin{fr}
\begin{ansblock}
It says nothing about where the minimum is. It describes the surface near the
point and has no opinion about anywhere else.
\end{ansblock}

Which is exactly Program~\ref{prog:P15}'s finding wearing different clothes:
the gradient is a local statement and the minimum's position is a global one.
The Hessian is local too, and adding it does not change that \dash{} it makes
the local description better, not wider.

\begin{rigourbox}
The statement that the error is $O(h^{3})$ needs the third derivative to exist
and be bounded near the point, which is Taylor's theorem with a remainder, and
it is in any first analysis course. This program uses the conclusion and
measures it; it does not prove it.

What is worth knowing without the proof is where it fails: at a point where
the function is not smooth, there is no quadratic model and the whole of this
program's argument has nothing to attach to. Program~\ref{prog:P16}'s $\relu$
at exactly zero is that case.
\end{rigourbox}
\end{fr}

\begin{fr}
Before the matrix, one question that has an answer you already know.

At a stationary point the gradient is zero, so the linear term vanishes and
the model is value plus curvature term alone. What does the sign of the
curvature tell you there?
\nextframe
\end{fr}

\begin{fr}
\begin{ansblock}
Positive is a minimum, negative a maximum \dash{} and in one dimension that
is the whole classification.
\end{ansblock}

Program~\ref{prog:F11} said a zero derivative is not a minimum and left the
test to \enquote{the derivative of the derivative}. This is that test, and in
one dimension it is complete.

In more than one it is not, and Program~\ref{prog:P10} already showed why:
there are shapes with no one-dimensional analogue. Getting from a sign to a
classification is what the next section's matrix is for.
\end{fr}

% ============================================================
\section{The matrix of second derivatives}
\label{sec:P17-hessian}
\index{Hessian}

\begin{fr}
With several inputs there is a second derivative for every \emph{pair} of
them: differentiate with respect to $x_i$, then with respect to $x_j$.

Arrange them as a matrix and it is the \textbf{Hessian}:
\[ H_{ij} = \frac{\partial^{2} f}{\partial x_i \, \partial x_j} \]

Program~\ref{prog:P16} said the derivative of a vector function is a matrix
with one row per output. Say what shape this one is, and why it is that shape.
\dotline
\nextframe
\end{fr}

\begin{fr}
\begin{ansblock}
Square, one row and one column per input \dash{} because it is the Jacobian of
the \emph{gradient}, which is a function from the inputs back to as many
outputs.
\end{ansblock}

So nothing new has been defined here either. The gradient takes $n$ inputs to
$n$ outputs, and Program~\ref{prog:P16}'s rule gives its derivative $n$ rows
and $n$ columns.

The second-order model in many dimensions is then the one-dimensional one with
the two products replaced by the two things
Program~\ref{prog:P05} and Program~\ref{prog:P10} supply:
\[ f(\mathbf{p} + \mathbf{h}) \approx f(\mathbf{p})
   + \nabla f \cdot \mathbf{h}
   + \tfrac{1}{2}\, \mathbf{h}^{\mathsf{T}} H \mathbf{h} \]
a dot product and a quadratic form.
\end{fr}

\begin{fr}
That quadratic form is the reason this program can talk about directions at
all, and it rests on one property of $H$. Which?
\nextframe
\end{fr}

\begin{fr}
\begin{ansblock}
It is \textbf{symmetric}: differentiating in the two orders gives the same
answer.
\end{ansblock}

And symmetry is what unlocks everything Program~\ref{prog:P10} proved. A
symmetric matrix has a full orthogonal set of eigenvectors and real
eigenvalues, so the surface near the point has a set of perpendicular
directions each with its own curvature, and the bowl-saddle-ridge
classification applies unchanged.

The script checks the two mixed partials agree to better than
$\val{p17.sym.bound}$ at a point of a function with no symmetry built into it.

\begin{note}
That check is corroboration and not the reason. Symmetry is a theorem
\dash{} it needs the second partials to be continuous near the point, and it
is in any first analysis course \dash{} and a numerical agreement at one point
would be poor evidence for a statement about every function.

The ceiling is committed rather than the measurement, for the reason
Program~\ref{prog:P06} paid for: the gap is rounding noise whose size is a
property of the machine, and here it happens to be exactly zero, which is
luck of this function rather than something another build owes you.
\end{note}
\end{fr}

\begin{fr}
So the classification Program~\ref{prog:P10} gave for quadratic forms is now a
statement about \emph{any} smooth function at a stationary point. Say it: what
do the signs of the Hessian's eigenvalues tell you?
\dotline
\nextframe
\end{fr}

\begin{fr}
\begin{ansblock}
All positive is a minimum, all negative a maximum, mixed signs a saddle, and a
zero eigenvalue means the second-order model does not decide.
\end{ansblock}

Which is Program~\ref{prog:P10}'s bowl, its inverted bowl, its saddle and its
ridge, in that order, and no new work was needed to get here \dash{} the
Hessian is the matrix that program was already talking about.

Now use it on the explanation everybody reaches for when a run stops
improving. A model has $\val{p17.params}$ parameters, so a stationary point of
its loss has that many eigenvalue signs. What has to be true of them for the
point to be a minimum, and in how many ways can that be missed by a single
sign?
\dotline
\nextframe
\end{fr}

\begin{fr}
\begin{ansblock}
Every single one has to be positive. It can be missed by one sign in
$\val{p17.params}$ ways, one per parameter \dash{} and those are only the
ways to miss it by one.
\end{ansblock}

\begin{trapbox}
\textbf{\enquote{My loss surface has local minima, and that is what training
gets stuck in.}} That is the two-dimensional picture carried somewhere it does
not survive the counting above: a minimum needs every sign to agree and a
saddle needs one dissenter, so a stationary point in high dimension is
overwhelmingly likely not to be a minimum at all.

What a run that has stopped improving is usually sitting in is one of two
other things. A saddle, which has a downhill direction and a gradient too
small to follow it quickly. Or a long narrow valley, which has a downhill
direction and is being crossed rather than followed \dash{} and that second
one is the next section's subject rather than any kind of minimum.

The honest limit of the argument: it is arithmetic about signs and not a
measurement of any trained network, and this book has not measured the latter.
What it does establish is that a zero gradient is much weaker evidence of a
minimum in a billion dimensions than in two.
\end{trapbox}
\end{fr}

% ============================================================
\section{Why the step size has a ceiling}
\label{sec:P17-stepsize}
\index{curvature!and step size}

\begin{fr}
Now the payoff, and Program~\ref{prog:P15} left it exactly one line away.

On a quadratic, gradient descent multiplies each eigendirection's coordinate
by $1 - \eta\lambda$ every step, with that direction's own $\lambda$. P15
measured $\val{p15.fac.hi}$ and $\val{p15.fac.lo}$ on its bowl and said
explicitly that where the boundary lies is this program's.

So find it. What does $1 - \eta\lambda$ have to satisfy for the coordinate not
to grow?
\dotline
\nextframe
\end{fr}

\begin{fr}
\begin{ansblock}
$\lvert 1 - \eta\lambda \rvert < 1$, which for positive $\eta$ and $\lambda$
is $\eta < \dfrac{2}{\lambda}$.
\end{ansblock}

\textbf{And that is the whole derivation.} Multiplying by a number repeatedly
grows without bound exactly when that number exceeds one in size; the number
here is $1 - \eta\lambda$; and rearranging gives the step size.

Nothing about neural networks, losses or training is in the argument. It is a
statement about a geometric sequence, which is Program~\ref{prog:F04}'s.

\mermaidfig{p17-why-the-step-caps}{One factor per direction, and what bounds
it.}{why the step caps}
\end{fr}

\begin{fr}
Now the part that makes it a statement about a whole surface rather than one
direction.

A bowl has several eigenvalues at once, and the walk has to stay bounded in
\emph{every} direction. On Program~\ref{prog:P10}'s bowl the two are
$\val{p17.lam.hi}$ and $\val{p17.lam.lo}$. What is the largest $\eta$ that
works, and which direction decided it?
\dotline
\nextframe
\end{fr}

\begin{fr}
\ans{$\val{p17.eta.star}$, set by the steepest}
Because each direction demands $\eta < 2/\lambda$ and the smallest of those
demands wins. The steep direction allows $\val{p17.eta.star}$; the shallow one
would allow $\val{p17.eta.shallow}$, which is $\val{p17.eta.gap}$ times
larger and is irrelevant, because taking it would blow the other one up.

\textbf{So the fastest-curving direction sets the step size for the whole
model, and every other direction lives with what it is given.} That is the
sentence Program~\ref{prog:P15}'s zig-zag was a picture of.

The script does not only derive it. It sweeps $\eta$ and finds the value at
which the walk first stops converging, which comes out at
$\val{p17.eta.measured}$ \dash{} the derived boundary, measured.
\end{fr}

\begin{fr}
And the gate, because two programs describing one walk must not disagree about
it. Program~\ref{prog:P15} walks that bowl at $\eta = \val{p17.p15.eta}$.
Against a boundary of $\val{p17.eta.star}$, is that walk safe?
\nextframe
\end{fr}

\begin{fr}
\begin{ansblock}
Yes, and only just \dash{} which is why its steep coordinate flips sign every
step instead of diverging.
\end{ansblock}

The script reads P15's committed step size and asserts it lies strictly under
this program's boundary, so the two cannot come apart.

\begin{trapbox}
\textbf{\enquote{The learning rate was too high} is usually said as a
judgement and it is an inequality.} What it means precisely is
$\eta > 2/\lambda_{\max}$, and the quantity on the right is a property of the
loss surface at the point you are standing on rather than of the optimiser or
the data.

Step one over the boundary and it is not a gentle degradation. Here is the
same bowl at $\eta = \val{p17.demo.eta}$, a hair above
$\val{p17.eta.star}$, for $\val{p17.demo.steps}$ steps:

\transcript{p17-too-big}

Each step multiplies by a number bigger than one in size, so the growth is
geometric and the sign alternates. \emph{Loss went to NaN} is that sequence
continuing until the format runs out, which Program~\ref{prog:P01} priced.
\end{trapbox}
\end{fr}

\begin{fr}
One more quantity falls out of the same two numbers, and it is one this book
has met twice already.

The ratio of the largest curvature to the smallest is
$\val{p17.kappa}$ here. What is that number called, and where has it appeared?
\dotline
\nextframe
\end{fr}

\begin{fr}
\begin{ansblock}
The \textbf{condition number}, of the Hessian \dash{} the same quantity
Program~\ref{prog:P11} defined for a matrix and Program~\ref{prog:P09} was
careful not to confuse with the determinant.
\end{ansblock}

There it measured how much a solve amplifies error. Here it measures how much
the step size you are allowed differs from the step size the slow direction
wants, and the two readings are the same arithmetic.

Its consequence is a rate. Even at the best compromise step
$\val{p17.best.eta}$, the slow direction shrinks by a factor of
$\val{p17.rate}$ per step, so closing $99$ per cent of the remaining gap takes
$\val{p17.steps.99}$ steps \dash{} on a bowl with only two directions and
exact arithmetic.

\textbf{That is the cost of a ravine stated as a number}, and it grows with the
condition number rather than with the size of the problem.
\end{fr}

% ============================================================
\section{Using the curvature, and what it costs}
\label{sec:P17-newton}
\index{Hessian!Newton's method}

\begin{fr}
If the trouble is that one step size has to serve directions with different
curvatures, there is an obvious fix: give each direction its own.

That is Newton's method \dash{} instead of stepping against the gradient,
solve $H\,\mathbf{s} = -\nabla f$ and step by $\mathbf{s}$. On a quadratic,
how many steps does it take?
\nextframe
\end{fr}

\begin{fr}
\ans{$\val{p17.newton.steps}$}
Exactly one, and it lands on the minimum. Dividing each direction by its own
curvature is precisely the correction the previous section said was missing,
and on a quadratic the model is the function, so the model's minimum is the
real one.

\textbf{So the method that solves the whole of section 3's problem exists and
is a page of first-year mathematics.} The question is only what it costs, and
the answer is why you have never used it.
\end{fr}

\begin{fr}
Count. A model has $\val{p17.params}$ parameters. How many entries does its
Hessian have, and what would that be in bytes at two bytes each?
\dotline
\nextframe
\end{fr}

\begin{fr}
\begin{ansblock}
$\val{p17.hess.entries}$ entries, so about $\val{p17.hess.bytes}$ bytes.
\end{ansblock}

Program~\ref{prog:F01}'s ladder of magnitudes has nothing on it that big. The
number of parameters is a number you can hold; the number of \emph{pairs} of
parameters is not, and squaring is what does it.

And storing it is the cheap part. Solving with it is cubic, about
$\val{p17.hess.solve}$ operations, which Program~\ref{prog:P09} priced in its
own terms: elimination builds $n$ numbers where the inverse builds $n^{2}$,
and here even $n$ is too large.

\textbf{So the Hessian is exactly like the Jacobian of
Program~\ref{prog:P16}: nobody forms one.} What is different is that for the
Jacobian the products people actually want are cheap, and for the Hessian the
thing people want \dash{} its inverse applied to the gradient \dash{} is not.
\end{fr}

\begin{fr}
\begin{note}
That is not the end of the subject, and it would be wrong to leave the
impression that second-order information is unusable. There is a family of
methods that never forms $H$ and builds an approximation to its action out of
the gradients it has already seen, and there are methods that use only the
diagonal.

Naming where the line falls is more useful than the list: \textbf{anything that
needs $n^{2}$ storage is out at this scale, and anything that needs $n$ is
in.} Program~\ref{prog:P20}'s optimisers all live on the right side of that
line, and the per-parameter scaling several of them do is the diagonal idea
under another name.
\end{note}
\end{fr}

\begin{fr}
There is a second objection to Newton's method that has nothing to do with
cost, and it matters more than people expect. Look again at what the step
does: it divides by the curvature in each direction.

What happens at a saddle?
\dotline
\nextframe
\end{fr}

\begin{fr}
\begin{ansblock}
Some curvatures are negative, so dividing by them steps \emph{towards} the
saddle in those directions rather than away.
\end{ansblock}

Newton's method finds stationary points, not minima, and it does not
distinguish between them \dash{} it is the multi-dimensional form of the fact
that Program~\ref{prog:F11} kept pointing at, that a zero derivative is not a
minimum.

\textbf{Combined with the counting argument in section 2}, that is a real
problem rather than a curiosity: in high dimension most stationary points are
saddles, so a method that converges to stationary points converges to the
wrong kind of thing. Every practical second-order method has to do something
about this, and what it does is a modification rather than the method itself.
\end{fr}

% ============================================================
\section{Sharp and flat, honestly}
\label{sec:P17-sharpness}
\index{curvature!sharp and flat minima}

\begin{fr}
Program~\ref{prog:P10} used the words \emph{ravine} and \emph{sharp minimum}
for basins with a large eigenvalue ratio. There is a well-known claim built on
those words: that flat minima generalise better than sharp ones.

This section is about which parts of that are statements about the trained
model. Start with the arithmetic. If you rescale one parameter \dash{} write
$w = c\,u$ and divide that parameter's input by $c$ \dash{} what happens to
the function the network computes?
\dotline
\nextframe
\end{fr}

\begin{fr}
\begin{ansblock}
Nothing. Every output is identical for every input, because the two factors of
$c$ cancel.
\end{ansblock}

The model is the same model. It is written in different coordinates, which is
Program~\ref{prog:P09}'s change of basis applied to a parameter rather than to
a vector.

Now the loss as a function of $u$ rather than of $w$. Its second derivative
picks up a factor of $c^{2}$ by the chain rule. So with $c = \val{p17.rescale.c}$
the curvature changes by what factor?
\nextframe
\end{fr}

\begin{fr}
\ans{$\val{p17.rescale.factor}$}
The script checks both halves: that the outputs agree to the last bit at every
input tried, and that the curvature is multiplied by exactly
$\val{p17.rescale.factor}$.

\textbf{So one model has two sharpnesses, and nothing distinguishes them
except how somebody chose to write the parameters.} A number that changes
while the function does not is a fact about the coordinates.

\mermaidfig{p17-sharp-is-not-a-fact}{The same function, two curvatures.}{sharp
is not a fact}
\end{fr}

\begin{fr}
\begin{trapbox}
\textbf{\enquote{This minimum is flatter, so it will generalise better}} is
not a claim you can evaluate until somebody says what \emph{flatter} is being
measured in.

A raw eigenvalue of the Hessian is not it, by the frame above. Nor is any
quantity built from the raw eigenvalues alone \dash{} the trace, the largest,
the ratio \dash{} because all of them move under a rescaling that changes
nothing about the model.

This is not a technicality that applies to unusual architectures. A network
with any elementwise scaling in it \dash{} which is nearly all of them
\dash{} admits exactly this reparameterisation, and normalisation layers make
it a symmetry the architecture has by construction.
\end{trapbox}
\end{fr}

\begin{fr}
The honest position is worth stating in full, because it is easy to over-
correct into saying the whole idea is empty.

What is established here is narrow: a bare sharpness number is not a property
of the function. That does not refute the empirical claim, and it does not
mean nothing survives.

What would it take for a sharpness-based claim to be about the model rather
than the coordinates?
\dotline
\nextframe
\end{fr}

\begin{fr}
\begin{ansblock}
A measure that transforms the same way the parameters do, so that rescaling
leaves it unchanged.
\end{ansblock}

Which is a real thing to ask for and a real thing to check, and it is the
question a careful paper on the subject answers explicitly and a careless one
does not mention.

\begin{warning}
\textbf{This book has not measured the generalisation claim and does not
adjudicate it.} Doing so needs trained models, a definition of flatness that
survives the frame above, and enough runs to beat variance \dash{} which is
Program~\ref{prog:P34}'s subject and not this one's.

What this program establishes is a filter, not a verdict: given a claim about
sharpness, ask whether the quantity it names would change under a rescaling
that leaves the model identical. If it would, the claim is about the
coordinates. That is a question you can answer in a minute and it disposes of
a large fraction of what gets said.
\end{warning}
\end{fr}

\begin{fr}
One last observation, and it closes the loop back to section 3.

The step-size bound $\eta < 2/\lambda_{\max}$ has the same problem, and
noticing that is the honest way to hold both.

Say why, and then say why it does not matter for that use.
\dotline
\nextframe
\end{fr}

\begin{fr}
\begin{ansblock}
Because $\lambda_{\max}$ rescales too. It does not matter because $\eta$
rescales with it, in the opposite direction, so the \emph{step actually taken}
is unchanged.
\end{ansblock}

That is the difference between the two uses in one sentence. In section 3 the
curvature appears alongside the step size, and the pair is invariant even
though neither member is. In section 5 the curvature is quoted alone, and
alone it means nothing.

\textbf{So the right question about any curvature number is never \enquote{is
it big} but \enquote{what is it big compared with}}, and section 3 has an
answer to that where the sharpness literature has to supply one.
\end{fr}

% ============================================================
\begin{summarybox}
\sumitem{1--2}{\textbf{The second-order Taylor model is value, slope and half
  the curvature times the square}, and the $\tfrac{1}{2}$ is what makes the
  model's second derivative match the function's.}
\sumitem{3--4}{\textbf{\enquote{Best} means an order, not a size}: the
  quadratic's error falls like $h^{\val{p17.order.quad}}$ where the line's
  falls like $h^{\val{p17.order.lin}}$, checked by halving $h$ rather than by
  quoting two errors.}
\sumitem{5--6}{\textbf{The model is local.} It makes the local description
  better, not wider, and says nothing about where the minimum is.}
\sumitem{7--8}{At a stationary point the linear term vanishes, so the sign of
  the curvature classifies \dash{} completely in one dimension, and not in
  more.}
\sumitem{9--10}{\textbf{The Hessian is the Jacobian of the gradient}, so
  Program~\ref{prog:P16}'s rule makes it square with one row and column per
  input.}
\sumitem{11--12}{\textbf{It is symmetric}, so Program~\ref{prog:P10} applies
  unchanged: real eigenvalues, perpendicular directions, each with its own
  curvature.}
\sumitem{13--14}{The eigenvalue signs give bowl, inverted bowl, saddle and
  ridge \dash{} P10's classification, now for any smooth function at a
  stationary point.}
\sumitem{15}{\textbf{A stationary point in high dimension is overwhelmingly
  unlikely to be a minimum}: a minimum needs every one of $\val{p17.params}$
  signs positive, and it can be missed by one sign in as many ways.
  \enquote{Stuck in a local minimum} is the two-dimensional picture carried
  somewhere it does not survive.}
\sumitem{16--17}{\textbf{The walk stays bounded exactly when
  $\lvert 1 - \eta\lambda\rvert < 1$, that is $\eta < 2/\lambda$}, which is a
  fact about a geometric sequence and about nothing else.}
\sumitem{18--19}{\textbf{The steepest direction sets the step for the whole
  model.} Here $\val{p17.eta.star}$ against the shallow direction's
  $\val{p17.eta.shallow}$, a factor of $\val{p17.eta.gap}$. Derived, then
  measured by sweeping: $\val{p17.eta.measured}$.}
\sumitem{20--21}{Program~\ref{prog:P15}'s walk at $\val{p17.p15.eta}$ is under
  the boundary, and the script asserts it. \textbf{\enquote{The learning rate
  was too high} is the inequality $\eta > 2/\lambda_{\max}$}, and past it the
  growth is geometric.}
\sumitem{22--23}{\textbf{The condition number $\val{p17.kappa}$ is the same
  quantity Program~\ref{prog:P11} defined}, read on the Hessian. At the best
  step it still takes $\val{p17.steps.99}$ steps to close $99$ per cent of the
  slow direction's gap.}
\sumitem{24--25}{\textbf{Newton's method lands on a quadratic's minimum in
  $\val{p17.newton.steps}$ step}, by giving each direction its own step size.}
\sumitem{26--28}{\textbf{And its Hessian has $\val{p17.hess.entries}$ entries
  at a real model size}, with a cubic solve on top. Nobody forms a Hessian
  either \dash{} and unlike the Jacobian, the product people want is not
  cheap.}
\sumitem{29--30}{\textbf{Newton finds stationary points, not minima}, so at a
  saddle it steps the wrong way in the negative directions. With section 2's
  counting argument that is a real problem in high dimension.}
\sumitem{31--33}{\textbf{Rescaling a parameter changes the curvature by
  $\val{p17.rescale.factor}$ and leaves every output identical.} So one model
  has two sharpnesses and a bare sharpness number is a fact about the
  coordinates.}
\sumitem{34}{\textbf{No quantity built from the raw eigenvalues survives it}
  \dash{} not the largest, the trace or the ratio \dash{} and normalisation
  layers make the rescaling a symmetry the architecture has by construction.}
\sumitem{35--36}{What would survive is a measure that transforms as the
  parameters do. \textbf{This book has not measured the generalisation claim
  and does not adjudicate it}; what it supplies is a filter.}
\sumitem{37--38}{\textbf{The step-size bound has the same problem and does not
  suffer from it}, because $\eta$ rescales with $\lambda$ and the step taken is
  unchanged. The question about a curvature is never \emph{is it big} but
  \emph{compared with what}.}
\end{summarybox}

\canyou

\begin{testexercises}
\item A quadratic has curvature $\lambda = 8$ in its only direction. Give the
  largest step size that converges, and say what happens at exactly that
  value.
  \answerto{$\eta < 2/8 = \num{0.25}$. At exactly $\num{0.25}$ the factor is
  $1 - 2 = -1$, so the coordinate flips sign and keeps its size forever
  \dash{} it neither converges nor diverges, which is the boundary being a
  boundary rather than a safety margin.}
\item A loss has Hessian eigenvalues $50$, $5$ and $\num{0.5}$ at the current
  point. Give the largest usable step, the condition number, and which
  direction is the problem.
  \answerto{$\eta < 2/50 = \num{0.04}$, set by the $50$. The condition number
  is $50/\num{0.5} = 100$. The problem direction is the $\num{0.5}$ one: it is
  allowed $\eta < 4$ and gets $\num{0.04}$, so it moves at a hundredth of the
  rate it could.}
\item Why does the derivation in section 3 never mention how many parameters
  the model has?
  \answerto{Because it is a statement about one eigendirection at a time, and
  the directions do not interact on a quadratic. Adding directions adds
  demands of the same form, and the binding one is still the largest
  curvature. The parameter count changes what $\lambda_{\max}$ is likely to
  be; it does not enter the argument.}
\item Somebody reports that their model converged to a minimum with maximum
  Hessian eigenvalue $\num{0.02}$, and concludes it is a flat minimum that
  will generalise well. Give the two questions you would ask.
  \answerto{First: in what coordinates? Rescaling a parameter moves that
  number by the square of the scale while the model is unchanged, so
  $\num{0.02}$ alone says nothing. Second: compared with what? A curvature is
  only meaningful against something with the same units, which in section 3
  is the step size. Neither question is hostile and both have answers; a claim
  that cannot answer them is about the coordinates.}
\item Your training run diverges. Give the arithmetic statement that
  corresponds to \enquote{the learning rate was too high}, and one thing other
  than lowering $\eta$ that would fix it.
  \answerto{$\eta > 2/\lambda_{\max}$ at the point the run was at. Other than
  lowering $\eta$: reduce $\lambda_{\max}$, which is what normalisation and a
  smaller initialisation both do, since the bound involves the surface as much
  as the optimiser. Note that $\lambda_{\max}$ is a property of where you are
  standing, so a run can be stable for a thousand steps and then not.}
\end{testexercises}

\begin{furtherproblems}
\item Show that on a quadratic with curvature $\lambda$, the step
  $\eta = 1/\lambda$ reaches the minimum in one move.
  \answerto{The factor is $1 - \eta\lambda = 0$, so the coordinate goes to
  zero in a single step. That is Newton's method in one dimension: dividing by
  the curvature is exactly choosing $\eta = 1/\lambda$, and the reason it does
  not generalise cheaply is that in many dimensions the division is a solve.}
\item Why is $\eta = 2/(\lambda_{\max} + \lambda_{\min})$ the best single step
  size for a quadratic, rather than something closer to $1/\lambda_{\max}$?
  \answerto{Because the worst direction is whichever of the two extremes has
  the larger $\lvert 1 - \eta\lambda\rvert$, and that worst case is minimised
  where the two are equal and opposite. Solving
  $1 - \eta\lambda_{\min} = \eta\lambda_{\max} - 1$ gives it. The resulting
  rate is $(\kappa - 1)/(\kappa + 1)$, which is where
  $\val{p17.rate}$ came from.}
\item The Hessian of a sum is the sum of the Hessians. What does that say
  about the curvature of a loss averaged over a batch?
  \answerto{That it is the average of the per-example curvatures, so the
  batch's $\lambda_{\max}$ is bounded by the largest per-example one and is
  usually much smaller, because the individual Hessians do not share their
  steep directions. That is one honest reason a larger batch tolerates a
  larger step \dash{} and it is a statement about the surface rather than
  about noise, which is Program~\ref{prog:P21}'s.}
\item Section 5's rescaling used one parameter. What does the same argument do
  to the whole Hessian if you rescale every parameter by $c$?
  \answerto{Multiplies it by $c^{2}$ throughout, so every eigenvalue scales
  and the condition number \dash{} a \emph{ratio} \dash{} does not. That is
  worth noticing: the condition number is invariant under a uniform rescaling
  where the individual eigenvalues are not, which is one reason it is the more
  quotable quantity. It is not invariant under rescaling parameters
  differently, which is the case section 5 built.}
\item Why does a zero eigenvalue of the Hessian leave the classification
  undecided, rather than meaning the direction is flat?
  \answerto{Because the second-order model says nothing there, and the
  function's behaviour in that direction is decided by terms the model
  dropped. $x^{4}$ and $-x^{4}$ both have zero second derivative at the
  origin and one has a minimum where the other has a maximum. That is the
  local model's limit showing itself, and it is why \enquote{flat direction}
  and \enquote{zero curvature} are not the same statement.}
\end{furtherproblems}
